\documentclass[11pt]{article}

\usepackage[margin=1in]{geometry}
\usepackage[T1]{fontenc}
\usepackage[utf8]{inputenc}
\usepackage{lmodern}
\usepackage{amsmath, amssymb, amsthm}
\usepackage{booktabs}
\usepackage{array}
\usepackage{enumitem}
\usepackage{hyperref}
\usepackage[table]{xcolor}

\hypersetup{
  colorlinks=true,
  linkcolor=blue,
  urlcolor=blue,
  citecolor=blue,
  pdftitle={HSA as Sparse Causal Routing and a Path Toward Typed Semantic Graphs},
  pdfauthor={OpenAI Codex}
}

\newcolumntype{C}{>{\centering\arraybackslash}m{1.5em}}
\newcommand{\mempty}{\cellcolor{gray!10}\rule{0pt}{1.5ex}\hspace{0.75em}}
\newcommand{\mlive}{\cellcolor{blue!20}\rule{0pt}{1.5ex}\hspace{0.75em}}
\newcommand{\mrowpick}{\cellcolor{green!25}\rule{0pt}{1.5ex}\hspace{0.75em}}
\newcommand{\mcolpick}{\cellcolor{violet!22}\rule{0pt}{1.5ex}\hspace{0.75em}}
\newcommand{\mcore}{\cellcolor{green!30}\rule{0pt}{1.5ex}\hspace{0.75em}}
\newcommand{\mtail}{\cellcolor{orange!30}\rule{0pt}{1.5ex}\hspace{0.75em}}

\title{HSA as Sparse Causal Routing\\
\large Formalization, Current Packing Strategy, and a Path Toward Typed Semantic Graphs}
\author{}
\date{April 2026}

\begin{document}
\maketitle

\begin{abstract}
This note formalizes the current hierarchical sparse attention (HSA) stack as a
decoder-causal sparse routing problem over a sequence graph, explains the
current dense HSA reference, schedule-based sparse realization, and packing
strategy, and distinguishes two operational HSA regimes that matter for current
benchmarking: a naive adjacency-driven sparse realization and a packed HSA path
that reuses precomputed cached sidecars. The note records the current isolated
speedups of packed HSA over naive HSA, places those results next to flat and
sliding-window FA4 baselines, and then sketches a realistic ``leapfrog''
extension path from binary hierarchical masks toward typed, directed, and
eventually signed semantic routing graphs that could support precursors to
causal reasoning in autoregressive language models. The central claim is that
the current implementation is not ``arbitrary graph attention,'' but it already
contains reusable pieces of a graph compiler, sparse runtime, and offline
sidecar system that can be repurposed for richer semantic graphs.
\end{abstract}

\section{Motivation}

The current HSA recipe uses semantic structure already present in text
representation---sentence anchors, section anchors, document anchors, or other
hierarchical boundary markers---to build a sparse causal attention graph. This
graph is then realized either as a dense semantic mask or as an equivalent
schedule-based sparse computation.

The long-range motivation is twofold:
\begin{enumerate}[leftmargin=1.5em]
    \item reduce attention cost while preserving relevance over long contexts,
    \item bias the model toward structured information flow rather than letting
    dense attention allocate probability mass over all causal predecessors.
\end{enumerate}

That is already useful for ``context rot'' mitigation, but it also opens a more
ambitious direction: moving from document hierarchy to typed semantic graphs and
eventually to typed directional message passing that can approximate parts of
causal reasoning without claiming true interventional semantics.

\section{Formalizing Current HSA}

\subsection{Sequence, graph, and support mask}

Let a decoder consume tokens $x_1, \dots, x_T$. Standard causal self-attention
defines a dense lower-triangular support:
\[
\mathcal{E}_{\mathrm{causal}} = \{(i,j) \mid 1 \le j \le i \le T\}.
\]

Current HSA replaces this with a sparse causal subgraph
\[
G = (V, E), \quad V = \{1, \dots, T\}, \quad E \subseteq \mathcal{E}_{\mathrm{causal}}.
\]

Equivalently, it defines a binary support mask
\[
S_{ij} =
\begin{cases}
1 & \text{if } (i,j) \in E,\\
0 & \text{otherwise.}
\end{cases}
\]

The graph $E$ is not learned online. It is compiled from token-side semantic
metadata such as \texttt{keep\_ids} and \texttt{hash\_ids}, which encode the
hierarchical structure of the document.

\subsection{Dense HSA reference}

Current dense HSA is a masked exact attention reference:
\[
\mathrm{score}_{ij} = \frac{q_i^\top k_j}{\sqrt{d}} + m_{ij},
\]
where
\[
m_{ij} =
\begin{cases}
0 & \text{if } S_{ij}=1,\\
-\infty & \text{if } S_{ij}=0.
\end{cases}
\]
Then
\[
\alpha_{ij} = \frac{\exp(\mathrm{score}_{ij})}
{\sum_{\ell : S_{i\ell}=1}\exp(\mathrm{score}_{i\ell})},
\qquad
o_i = \sum_{j : S_{ij}=1} \alpha_{ij} v_j.
\]

This is the clean semantic reference: all questions about correctness reduce to
whether the dense mask implements the intended graph.

\subsection{Schedule-based sparse HSA}

The schedule-based sparse implementation is supposed to compute exactly the same
function as the dense mask, but using explicit sparse descriptors rather than a
dense masked traversal. In the repository, that compiler/runtime split lives
primarily in:
\begin{itemize}[leftmargin=1.5em]
    \item \texttt{flash\_attn/cute/hsa.py},
    \item \texttt{flash\_attn/cute/hsa\_cached\_2d\_forward\_analysis.py},
    \item \texttt{nanochat/hsa\_batch\_cache.py}.
\end{itemize}

The schedule is therefore a \emph{compiled sparse realization of the same
binary support graph}, not a different model class.

\section{What the Current Graph Semantics Actually Encode}

Current HSA is not general graph attention. It is a decoder-causal sparse graph
with a specific information-routing policy derived from hierarchy.

In the present recipe, the graph includes:
\begin{itemize}[leftmargin=1.5em]
    \item same-sentence body-level connectivity,
    \item same-section anchor connectivity,
    \item same-document anchor connectivity,
    \item cross-level edges such as body $\rightarrow$ previous sentence anchor,
    \item body $\rightarrow$ previous section anchor,
    \item sentence anchor $\rightarrow$ previous section anchor.
\end{itemize}

These cross-level edges matter. Without them, the graph is too lossy: a purely
same-level tree-shaped routing policy under-connects the model and discards
useful long-range semantic pathways.

The important subtlety is that this graph is a \emph{routing graph}, not an
ontology. The document hierarchy is a tree. Useful attention flow is closer to a
causalized directed acyclic graph over the prefix.

\section{Dense HSA Versus Sparse HSA}

\subsection{Dense HSA}

Dense HSA means: build the exact semantic mask, run the ordinary dense masked
attention reference, and let the kernel traverse all entries in the lower
triangle even if many entries are masked out. This is conceptually simple and is
the best semantic oracle when debugging graph meaning.

\subsection{Sparse HSA}

Sparse HSA means: compile the same semantic mask into sparse descriptors and
realize the same attention via sparse traversal, sparse runtime metadata, or
packed subproblems. Sparse HSA is therefore a systems optimization problem on
top of the same graph semantics.

\subsection{Invariant that matters}

The critical invariant is:
\[
\text{dense semantic mask} \equiv \text{sparse schedule} \equiv \text{kernel support}.
\]

When these differ, training behavior diverges even if the intended high-level
semantics were correct. Much of the engineering difficulty in HSA is therefore
not inventing a graph, but proving that every layer of the implementation
realizes the same graph.

\section{Why Packing Helps Forward More Than Backward}

\subsection{Forward is naturally query-owned}

Forward attention on a compressed subproblem is naturally row-owned:
\[
Q_b \in \mathbb{R}^{r \times d}, \quad
K_b, V_b \in \mathbb{R}^{u \times d},
\]
with one output $O_b$ and one log-sum-exp state per query row. Packing reduces
an irregular sparse neighborhood into a regular dense microproblem.

That is why forward benefits strongly from packing:
\begin{itemize}[leftmargin=1.5em]
    \item gather and pack once,
    \item compute one tiny dense exact attention,
    \item write one output and one LSE per row.
\end{itemize}

\subsection{Backward has an ownership mismatch}

Backward is harder because the natural ownership splits:
\[
dQ = dS K \quad \text{(query-row owned)}
\]
but
\[
dK \mathrel{+}= dS^\top Q, \qquad
dV \mathrel{+}= P^\top dO
\quad \text{(key-row owned)}.
\]

This means the same forward packing that makes a tiny dense forward problem does
not automatically make a tiny dense backward problem with a natural writeback
owner. The moment we leave the query-owned world, we pay for reductions,
atomics, or extra saved-state movement.

\subsection{Why \texttt{sparse\_mask} backward is still so good}

Even though \texttt{sparse\_mask} backward is coarser and often less elegant
than the most compressed forward path, it remains strong for several reasons:
\begin{enumerate}[leftmargin=1.5em]
    \item it stays close to the exact support graph;
    \item it amortizes work in large regular backward kernels instead of carrying
    fragile packed forward state into backward;
    \item it avoids the q-owned / k-owned mismatch becoming a bespoke packed
    writeback problem;
    \item once numerical issues are fixed, it is robust and easy to validate
    against the dense masked reference;
    \item empirically, the forward gains from aggressive packing are large,
    while the backward gains from the same packing are usually small or negative.
\end{enumerate}

Historically this is why the project leaned on compressed forward plus a more
conservative backward. In the current codebase, however, the practical
long-context training configuration is no longer ``plain
\texttt{sparse\_mask}'' as such; it is the packed HSA path that reuses cached
forward sidecars and the dedicated long-context backward route attached to that
path.

\subsection{Long-context caveat: the generic block-sparse \texorpdfstring{$dQ$}{dQ} path was the real bug}

The long-context debugging pass went through two phases. The first phase
correctly established that the dense \texttt{mask\_mod} fallback backward is too
slow for real long-context training. The second phase refined that diagnosis:
the problem was not ``fast backward'' in the abstract, but the \emph{generic}
FA4 block-sparse $dQ$ path used by the exact packed-mask traversal.

The key measurement trap was that the original isolated \texttt{fwdbwd}
benchmark was profiling a degenerate backward regime. At step $0$, the
attention output projection \texttt{c\_proj} is zero-initialized, so the
upstream gradient flowing into the attention output is effectively zero. Because
the implementation has an exact zero-\texttt{dout} fast path, that benchmark
mostly measured:
\begin{itemize}[leftmargin=1.5em]
    \item the real packed forward cost, plus
    \item an almost trivial backward.
\end{itemize}

Once we added a profiler mode that performs a real optimizer update before the
profiled step, the true long-context picture became obvious. The dense-traversal
fallback backward was indeed safe but far too slow, but the first fast
replacement---the generic FA4 block-sparse backward used by
\texttt{\_run\_hsa\_packed\_mask\_backward(...)}---was still wrong.

The decisive tensor-level probe was a direct comparison on the exact same
post-update $32768$ tensors:
\begin{itemize}[leftmargin=1.5em]
    \item $dK$ and $dV$ from the generic block-sparse traversal matched the
    dense-traversal oracle to numerical noise;
    \item $dQ$ did not.
\end{itemize}
Typical differences were
\[
\max|dQ_{\mathrm{dense}} - dQ_{\mathrm{sparse}}| \approx 5.5 \text{ to } 7.7,
\qquad
\mathrm{mean}|dQ_{\mathrm{dense}} - dQ_{\mathrm{sparse}}| \approx 0.15 \text{ to } 0.18,
\]
even though the upstream \texttt{dout} on that call was only of order
$10^{-8}$. That is not ordinary BF16 drift; it points to a wrong or stale
$dQ$ accumulation path.

\subsection{What actually worked}

The important refinement is that the \emph{dedicated HSA long-context backward}
paths already in the repository were not affected in the same way.

\begin{itemize}[leftmargin=1.5em]
    \item Sparse exact HSA with \texttt{legacy\_packed} mode uses the dedicated
    HSA long-context backward implementation.
    \item Cached generalized forward can now route its backward through the
    dedicated packed HSA backward as well.
\end{itemize}

On a real Gutenberg cache at $T=32768$, first real post-update step:
\begin{itemize}[leftmargin=1.5em]
    \item correctness-safe dense-traversal sparse path:
    \[
    \text{loss}_1 = 6.4592276,\qquad \text{step}_1 \approx 22.2\ \mathrm{s}
    \]
    \item dedicated HSA \texttt{legacy\_packed} sparse path:
    \[
    \text{loss}_1 = 6.4592276,\qquad \text{step}_1 \approx 10.2\ \mathrm{s}
    \]
\end{itemize}
and the observed gradient drift stayed at BF16 scale rather than exploding in
$dQ$.

On a fixed-seed resident five-step run at the same length:
\begin{itemize}[leftmargin=1.5em]
    \item correctness-safe sparse:
    final loss $\approx 3.32278$,
    throughput $\approx 1.37\mathrm{k}$ tok/s;
    \item \texttt{legacy\_packed} sparse:
    final loss $\approx 3.32252$,
    throughput $\approx 6.64\mathrm{k}$ tok/s;
    \item cached generalized forward plus dedicated packed HSA backward:
    final loss $\approx 3.32252$,
    throughput $\approx 8.81\mathrm{k}$ tok/s.
\end{itemize}

\subsection{Current policy}

The practical policy is therefore no longer ``dense fallback everywhere until
the long-context bug is fixed.'' It is:
\begin{enumerate}[leftmargin=1.5em]
    \item keep the generic FA4 block-sparse backward opt-in only, because its
    $dQ$ path is still not trustworthy at long context;
    \item keep short contexts on the simple correctness-first path;
    \item default long contexts to the dedicated HSA \texttt{legacy\_packed}
    backward once sequence length crosses the auto threshold;
    \item let cached generalized forward reuse that same dedicated packed HSA
    backward so the packed forward win survives end-to-end.
\end{enumerate}

The current auto threshold is controlled by
\texttt{FLASH\_ATTN\_HSA\_AUTO\_LEGACY\_PACKED\_BWD\_MIN\_SEQLEN} and defaults
to $32768$.

\subsection{What changed in the implementation}

The practical speedup came from routing changes, not from changing the HSA
algebra. There are now three relevant backward families:
\begin{itemize}[leftmargin=1.5em]
    \item dense traversal over the exact packed mask modifier;
    \item the generic FA4 block-sparse traversal driven by
    \texttt{block\_sparse\_tensors};
    \item the dedicated HSA long-context backward implementations
    (\texttt{legacy\_packed} and related HSA-specific paths).
\end{itemize}

The exact packed-mask helper still retains the conservative
\texttt{block\_sparse\_tensors=None} route as the dense oracle, and the generic
block-sparse traversal remains available for targeted kernel debugging. The
important new routing change is that long-context HSA training no longer has to
choose between ``slow but correct dense traversal'' and ``fast but wrong generic
block-sparse $dQ$.'' The repository now routes long contexts toward the
dedicated HSA packed backward by default, while keeping the generic FA4
block-sparse path explicit and opt-in.

\section{Row Packing, Column Packing, and Generalized 2-D Packing}

\subsection{Row packing}

Row packing compresses multiple query rows with similar neighborhoods into a
single regular microproblem. Informally:
\begin{itemize}[leftmargin=1.5em]
    \item choose a small set of query rows,
    \item take the union of their live keys,
    \item gather that union once,
    \item run one tiny dense attention over the packed rows.
\end{itemize}

This is the most natural packing because forward attention is query-owned.

\subsection{Column packing}

Column packing focuses on the key/value side:
\begin{itemize}[leftmargin=1.5em]
    \item compress repeated or shared key neighborhoods,
    \item reduce redundant K/V gathers,
    \item try to reuse union-K structure across rows.
\end{itemize}

This helps when many rows share similar key support, but on its own it is less
complete than row packing because it does not resolve the query-side ownership
pattern.

\subsection{Generalized 2-D packing}

Generalized 2-D packing treats the support matrix as a sparse binary matrix and
tries to extract dense submatrices:
\begin{itemize}[leftmargin=1.5em]
    \item exact dense tiles when a row group shares a common live key set,
    \item tiled residual tails for uncovered support,
    \item range-level fused execution where exact tiles and tails share one
    online softmax state.
\end{itemize}

That is more general than either pure row packing or pure column packing:
\begin{itemize}[leftmargin=1.5em]
    \item row packing asks ``which rows should be grouped?''
    \item column packing asks ``which K/V neighborhoods should be reused?''
    \item generalized 2-D packing asks ``which dense submatrices can be
    extracted from the sparse support matrix, and how should the remainder be
    tiled?''
\end{itemize}

The current cached generalized forward path in the repository is exactly this
kind of 2-D extraction:
\begin{itemize}[leftmargin=1.5em]
    \item exact dense core tiles,
    \item fused tiled tail,
    \item one range-local softmax state,
    \item one output write per row range.
\end{itemize}

\subsection{What packed forward actually computes}

The easiest way to say it precisely is: packing does \emph{not} change the
attention rule. It changes the layout of the live support so that dead
``white-space'' entries are physically removed before the kernel runs.

Take one sparse support subproblem with query rows $R=\{r_1,\dots,r_m\}$ and
key rows $C=\{c_1,\dots,c_n\}$. Packing defines two gather maps
\[
g_Q(a)=r_a,\qquad g_K(b)=c_b,
\]
then gathers
\[
\widetilde Q_a = Q_{g_Q(a)},\qquad
\widetilde K_b = K_{g_K(b)},\qquad
\widetilde V_b = V_{g_K(b)}.
\]
It also defines the packed binary mask
\[
\widetilde M_{ab} =
\begin{cases}
1 & \text{if } (g_Q(a),g_K(b)) \text{ is live in the original HSA graph},\\
0 & \text{otherwise.}
\end{cases}
\]

Then the packed microproblem is just exact masked attention on the reindexed
live rows and columns:
\[
\widetilde S_{ab} = \frac{\widetilde Q_a^\top \widetilde K_b}{\sqrt d},
\qquad
\widetilde P_{ab}
=
\frac{\exp(\widetilde S_{ab})\widetilde M_{ab}}
{\sum_c \exp(\widetilde S_{ac})\widetilde M_{ac}},
\qquad
\widetilde O_a = \sum_b \widetilde P_{ab}\widetilde V_b.
\]

So the intuitive picture
\begin{quote}
identify sparse tiles, remove the white space, fill the survivors into a dense
tile, run GEMM, and then reindex so the softmax is correct
\end{quote}
is basically right, with one important refinement: the softmax is not
``corrected afterward'' in an ad hoc way. The packed row/column indices and the
packed mask \(\widetilde M\) define the exact local normalization domain. When a
single original HSA row is split across several packed families or ranges, the
implementation carries an exact online softmax state (LSE / running max-sum)
across those pieces and merges them before the final writeback. That is what
keeps the packed forward algebraically identical to the original sparse row.

\subsection{What packed backward actually computes}

Backward uses the same gather maps, but the ownership is no longer as clean as
in forward. In packed coordinates,
\[
d\widetilde Q = d\widetilde S\,\widetilde K,\qquad
d\widetilde K \mathrel{+}= d\widetilde S^\top \widetilde Q,\qquad
d\widetilde V \mathrel{+}= \widetilde P^\top d\widetilde O.
\]
The writeback is then:
\begin{itemize}[leftmargin=1.5em]
    \item scatter \(d\widetilde Q\) back to the original query rows \(g_Q(a)\);
    \item scatter-add \(d\widetilde K\) and \(d\widetilde V\) back to the
    original key rows \(g_K(b)\).
\end{itemize}

This is the core asymmetry:
\begin{itemize}[leftmargin=1.5em]
    \item forward is naturally query-owned, so packing mostly gives a small
    exact dense microproblem plus one output write per row;
    \item backward must reduce many packed contributions back into shared
    original \(K/V\) rows, so the host-side descriptor path, batching policy,
    and scatter-add structure matter much more.
\end{itemize}

That is why ``packing helps forward'' does not automatically imply ``the same
packing gives an equally clean backward.'' The algebra is exact in both cases;
the systems problem is harder in backward because the packed representation has
to support cheap accumulation back into the original coordinates.

\subsection{Packing illustrations}

Figure~\ref{fig:packing-illustrations} shows the three packing viewpoints on a
small support matrix. Blue cells are live edges in the original sparse support.
Green cells indicate an extracted dense core, purple cells indicate a reused
column family, and orange cells indicate a tiled residual tail.

\begin{figure}[ht]
\centering
\setlength{\tabcolsep}{2pt}
\renewcommand{\arraystretch}{1.1}

\begin{minipage}[t]{0.31\textwidth}
\centering
\textbf{Row packing}\\[0.35em]
\begin{tabular}{c|*{6}{C}}
 & $k_1$ & $k_2$ & $k_3$ & $k_4$ & $k_5$ & $k_6$ \\
\hline
$q_1$ & \mrowpick & \mrowpick & \mempty & \mrowpick & \mempty & \mempty \\
$q_2$ & \mrowpick & \mrowpick & \mempty & \mrowpick & \mempty & \mempty \\
$q_3$ & \mempty & \mempty & \mlive & \mempty & \mlive & \mempty \\
$q_4$ & \mempty & \mempty & \mlive & \mempty & \mlive & \mlive \\
\end{tabular}

\vspace{0.5em}

\begin{tabular}{c|*{3}{C}}
 & $u_1$ & $u_2$ & $u_3$ \\
\hline
$q_1$ & \mcore & \mcore & \mcore \\
$q_2$ & \mcore & \mcore & \mcore \\
\end{tabular}

\vspace{0.35em}
\raggedright
\footnotesize
Select similar rows $\{q_1,q_2\}$, form the union
$\{k_1,k_2,k_4\}$ once, and run one dense $2 \times 3$ microproblem.
\end{minipage}
\hfill
\begin{minipage}[t]{0.31\textwidth}
\centering
\textbf{Column packing}\\[0.35em]
\begin{tabular}{c|*{6}{C}}
 & $k_1$ & $k_2$ & $k_3$ & $k_4$ & $k_5$ & $k_6$ \\
\hline
$q_1$ & \mcolpick & \mcolpick & \mempty & \mcolpick & \mempty & \mempty \\
$q_2$ & \mcolpick & \mcolpick & \mempty & \mcolpick & \mempty & \mempty \\
$q_3$ & \mempty & \mempty & \mlive & \mempty & \mlive & \mempty \\
$q_4$ & \mempty & \mempty & \mlive & \mempty & \mlive & \mlive \\
\end{tabular}

\vspace{0.5em}

\begin{tabular}{c|*{3}{C}}
 & $k_1$ & $k_2$ & $k_4$ \\
\hline
reuse & \mcolpick & \mcolpick & \mcolpick \\
\end{tabular}

\vspace{0.35em}
\raggedright
\footnotesize
Identify a repeated K/V family and gather those columns once. This reduces
redundant K/V movement, but row ownership is still unresolved by itself.
\end{minipage}
\hfill
\begin{minipage}[t]{0.31\textwidth}
\centering
\textbf{Generalized 2-D packing}\\[0.35em]
\begin{tabular}{c|*{6}{C}}
 & $k_1$ & $k_2$ & $k_3$ & $k_4$ & $k_5$ & $k_6$ \\
\hline
$q_1$ & \mcore & \mcore & \mempty & \mcore & \mempty & \mempty \\
$q_2$ & \mcore & \mcore & \mempty & \mcore & \mempty & \mempty \\
$q_3$ & \mempty & \mempty & \mtail & \mempty & \mtail & \mempty \\
$q_4$ & \mempty & \mempty & \mtail & \mempty & \mtail & \mtail \\
\end{tabular}

\vspace{0.5em}

\begin{tabular}{c|*{3}{C}}
 & $u_1$ & $u_2$ & $u_3$ \\
\hline
core & \mcore & \mcore & \mcore \\
tail & \mtail & \mtail & \mtail \\
\end{tabular}

\vspace{0.35em}
\raggedright
\footnotesize
Extract dense submatrices from the support itself: an exact dense core tile for
$(q_1,q_2)$, then cover the remainder with tiled residual tails that share one
range-local softmax state.
\end{minipage}

\caption{Three views of packing. In every case the blue/green/orange cells are
\emph{live} support and the white cells are dead support that is physically
removed from the packed microproblem. Row packing groups similar query rows,
column packing reuses repeated K/V families, and generalized 2-D packing
extracts dense submatrices from the support matrix and covers the remainder
with fused tail tiles. The packed microproblem is still exact masked attention
on the reindexed live rows and columns, and online softmax/LSE merging keeps
normalization identical to the original sparse row even when one row is split
across several packed ranges.}
\label{fig:packing-illustrations}
\end{figure}

The important systems point is that all three methods operate on the support
graph rather than on the exact scoring rule. This means the graph compiler,
sidecars, schedule families, and most of the current packing machinery remain
reusable if the score becomes relation-aware or typed. The part that changes is
the kernel math inside each packed family: score evaluation, value transforms,
and possibly normalization if the model moves beyond plain masked softmax.

\section{Benchmark Baselines and Preliminary Results}

\subsection{Benchmark labels}

The current benchmark story is easiest to understand if the runtime modes are
named by \emph{what work is actually being done}:
\begin{itemize}[leftmargin=1.5em]
    \item \textbf{flat FA4}: ordinary dense causal FA4 with no hierarchical
    sparsity;
    \item \textbf{sliding\_log}: a causal sliding-window FA4 baseline with
    window width $\lceil \log_2 T \rceil$;
    \item \textbf{strict\_log}: the same logarithmic sliding-window width
    $\lceil \log_2 T \rceil$ applied at \emph{every} layer, including the final
    layer;
    \item \textbf{sliding\_fixed}: a causal sliding-window FA4 baseline with a
    fixed window width of $1000$ tokens;
    \item \textbf{HSA naive}: exact HSA with cached-forward fallback disabled,
    i.e.\ the adjacency-driven sparse realization;
    \item \textbf{HSA packed}: exact HSA with cached-forward fallback enabled,
    i.e.\ the packed HSA path that reuses precomputed forward sidecars and the
    corresponding long-context backward route.
\end{itemize}

This distinction matters because once runtime fallback to precomputed cached
sidecars is enabled, a plain \texttt{sparse\_mask} benchmark is no longer a
clean ``naive sparse'' baseline on cache-backed batches. In current cache-backed
runs, the default exact HSA path can already reuse the packed cached forward
representation. For that reason, the meaningful systems comparison is now
\textbf{HSA naive} versus \textbf{HSA packed}, not \texttt{sparse\_mask} versus
an older explicit cached preset.

The distinction between \textbf{sliding\_log} and \textbf{strict\_log} is
operationally important. \textbf{sliding\_log} is not a pure local baseline; it
is a logarithmic local stack plus one dense global final layer. By contrast,
\textbf{strict\_log} is a pure log-band local stack. Recent Gutenberg, FineWeb,
and Wikipedia slices suggest that both are much stronger than a toy sparse
strawman and should remain in the benchmark set.

\subsection{Isolated packed-versus-naive HSA}

Table~\ref{tab:isolated-packed-vs-naive} gives the clean isolated comparison on
cache-backed resident batches. These are the same HSA semantics, same cached
batches, same model, and same losses to numerical noise; the only difference is
whether the runtime is forced onto the naive adjacency-driven sparse path or is
allowed to consume the packed cached sidecars.

\begin{table}[ht]
\centering
\begin{tabular}{lrrrrr}
\toprule
$T$ & naive dt (ms) & packed dt (ms) & packed tok/s & speedup & loss match \\
\midrule
$65536$ & $6549.4$ & $118.2$ & $554.6$k & $55.4\times$ & yes \\
$131072$ & $15657.0$ & $207.5$ & $631.6$k & $75.4\times$ & yes \\
$262144$ & $29184.6$ & $344.7$ & $760.6$k & $84.7\times$ & yes \\
\bottomrule
\end{tabular}
\caption{Isolated HSA systems comparison on cache-backed resident batches. The
packed path is the intended method; the naive path is the adjacency-driven
reference that does not reuse packed cached sidecars.}
\label{tab:isolated-packed-vs-naive}
\end{table}

The key result is not subtle: the packed path is already doing the intended
systems job. The speedup over the naive exact sparse realization grows with
sequence length, from roughly $55\times$ at $64$K to roughly $85\times$ at
$262144$, while preserving the same hierarchical support semantics.

\subsection{True multi-batch resident training runs}

The earlier repeated-single-batch resident smokes were useful for quick systems
checks, but they were too degenerate to say much about relative learning
behavior. Table~\ref{tab:true-multibatch-resident} replaces that with
\emph{true multi-batch} resident training on a scaled Gutenberg corpus
(\texttt{gutenberg\_large}, $31$ books, $\approx 23.6$M characters, and real
long-context train-cache budgets of $227$ batches at $64$K, $116$ at $128$K,
and $59$ at $262144$). To keep the resident benchmark memory-bounded while
still avoiding repeated-single-batch memorization, the runs below cycle over
selected resident subsets:
\begin{itemize}[leftmargin=1.5em]
    \item $64$K: $8$ distinct resident train batches;
    \item $128$K: $3$ distinct resident train batches;
    \item $262144$: $2$ distinct resident train batches.
\end{itemize}
All rows below use real cache-backed batches and cycle over multiple distinct
resident schedules rather than memorizing a single cached example.

\begin{table}[ht]
\centering
\begin{tabular}{llrrrr}
\toprule
$T$ & mode & train batches & final loss & mean dt (ms) & tok/s \\
\midrule
$65536$ & flat FA4 & $8$ & $2.9195$ & $299.7$ & $218.7$k \\
$65536$ & sliding\_log & $8$ & $2.3961$ & $65.4$ & $1001.8$k \\
$65536$ & sliding\_fixed & $8$ & $2.8806$ & $60.3$ & $1087.0$k \\
$65536$ & HSA packed & $8$ & $2.4858$ & $130.2$ & $503.3$k \\
\midrule
$131072$ & flat FA4 & $3$ & $2.8504$ & $947.2$ & $138.4$k \\
$131072$ & sliding\_log & $3$ & $1.7216$ & $199.6$ & $656.7$k \\
$131072$ & sliding\_fixed & $3$ & $2.8088$ & $306.2$ & $428.1$k \\
$131072$ & HSA packed & $3$ & $1.8358$ & $178.8$ & $733.2$k \\
\midrule
$262144$ & flat FA4 & $2$ & $2.8278$ & $1740.8$ & $150.6$k \\
$262144$ & sliding\_log & $2$ & $1.3794$ & $358.4$ & $731.3$k \\
$262144$ & sliding\_fixed & $2$ & $2.6345$ & $367.8$ & $712.7$k \\
$262144$ & HSA packed & $2$ & $1.4889$ & $351.6$ & $745.6$k \\
\bottomrule
\end{tabular}
\caption{True multi-batch resident training runs on the scaled
\texttt{gutenberg\_large} long-context caches. These are still internal
benchmark runs rather than full downstream training jobs, but they are
materially stronger than repeated-single-batch smokes because they cycle over
multiple distinct resident schedules from a larger real corpus.}
\label{tab:true-multibatch-resident}
\end{table}

Several points matter immediately.
\begin{enumerate}[leftmargin=1.5em]
    \item On the scaled corpus, the packed HSA path is already clearly ahead of
    dense flat FA4 by $64$K on both throughput and loss, and by $128$K and
    $262144$ the throughput gap is large: about $5.3\times$ at $128$K and
    about $5.0\times$ at $262144$, while also landing at materially lower loss
    on these short internal runs.
    \item The sliding-window baselines are much stronger than a toy strawman.
    They are not hierarchical and they should eventually lose on relevance when
    the task genuinely requires nonlocal graph structure, but they are very fast
    and must be included in any honest benchmark table.
    \item On these short true multi-batch resident runs, \texttt{sliding\_log}
    remains the strongest quality baseline. That is a useful warning: any claim
    that HSA ``obviously'' dominates sliding windows needs to be supported by
    longer and less degenerate training, not by intuition alone.
    \item The earlier scaled-corpus $64$K anomaly was not a missing packed
    dispatch. It came from a bad interaction between first-use per-schedule
    costs and a subset of packed sidecars that chose grouped
    \texttt{k\_window}/direct families instead of the fast \texttt{union\_2d}
    family. Rebuilding the selected $64$K sidecars with a benchmark-only
    policy override that forces the efficient unionized family, and warming all
    selected resident batches before timing, restores the expected packed-HSA
    regime.
\end{enumerate}

The $262144$ row remains the clearest systems crossover. Packed HSA cleanly
clears the dense flat baseline on both throughput and loss while landing
slightly ahead of both sliding baselines on throughput, though still behind
\texttt{sliding\_log} on loss. That still does not prove a general downstream
quality win by itself, but it is the right systems crossover: the packed
hierarchical path is now winning in the regime it was designed for, and the
remaining work is to run longer multi-batch training on larger corpora and to
improve smaller long-context regimes without relying on benchmark-specific
sidecar retuning.

\subsection{Issues encountered after the Gutenberg crossover}

The Gutenberg table established that packed HSA can already cross over dense
flat FA4 in the long-context regime. The next debugging pass on Wikipedia and
FineWeb exposed a different issue: the packed representation is
dataset-geometry-sensitive.

The main problems were:
\begin{enumerate}[leftmargin=1.5em]
    \item some wiki/web batches produced packed sidecars with large grouped
    \texttt{k\_window} / direct families rather than the compact
    \texttt{union\_2d} regime that worked well on Gutenberg;
    \item those sidecars often lacked a cheap cached fused-backward route, so
    the runtime fell back to the legacy packed backward machinery;
    \item the remaining runtime bottleneck on those batches was not the dense
    GEMM itself, but host-side descriptor building and thousands of tiny packed
    backward panel calls.
\end{enumerate}

One concrete improvement already landed: sorting and coalescing legacy packed
backward panels by shape before execution. On a representative Wikipedia $64$K
batch, that reduced the hybrid backward batch count from roughly
\[
165 \;\longrightarrow\; 15
\]
under the current default merge budget, and reduced the first real post-update
step from about
\[
9.17\ \mathrm{s} \;\longrightarrow\; 1.65\ \mathrm{s}.
\]

That is an important refinement of the current state:
\begin{itemize}[leftmargin=1.5em]
    \item the remaining non-Gutenberg blocker is not ``HSA semantics are
    wrong'';
    \item it is also not ``GEMM packing is inherently slower than FA4'';
    \item it is a systems issue in the packed backward descriptor/batching path
    on less favorable support geometries.
\end{itemize}

So the current non-Gutenberg work is about making the packed backward path less
fragile to wiki/web geometry, not about changing the HSA graph definition.

\subsection{MFU caveat}

The benchmark harness currently reports MFU using a dense-equivalent FLOP
estimate. That is serviceable for flat and sliding-window FA4 rows, but it is
not publishable for HSA rows because the packed and sparse paths do not execute
the dense-equivalent workload. For now, the trustworthy systems metrics are:
\begin{itemize}[leftmargin=1.5em]
    \item mean step time,
    \item tokens per second,
    \item peak memory,
    \item and loss trajectories.
\end{itemize}

\section{What the Current Implementation Already Gives Us}

The current stack contains several reusable components that are more general
than ``hierarchy only'' suggests.

\subsection{1. A graph compiler from token annotations to sparse causal support}

Today this compiler is expressed through:
\begin{itemize}[leftmargin=1.5em]
    \item \texttt{keep\_ids},
    \item \texttt{hash\_ids},
    \item \texttt{compute\_hsa\_mask(...)},
    \item \texttt{build\_hsa\_schedule(...)}.
\end{itemize}

This is already a graph compiler. It just currently compiles one specific
family of hierarchical binary support graphs.

\subsection{2. Exact dense and sparse references}

The repository already has the conceptual machinery needed to separate:
\begin{itemize}[leftmargin=1.5em]
    \item intended semantics,
    \item sparse schedule,
    \item kernel behavior.
\end{itemize}

That is crucial for any richer graph scheme. If the graph is extended to typed
or directed edges, the same discipline will still be necessary.

\subsection{3. Sidecars and offline precomputation}

The schedule sidecar and forward sidecar system is not specific to hierarchy.
It is a reusable vehicle for:
\begin{itemize}[leftmargin=1.5em]
    \item offline graph compilation,
    \item offline packing,
    \item runtime attachment of structural metadata,
    \item avoiding host-side rebuild costs.
\end{itemize}

This becomes even more valuable if the graph moves from simple hierarchy to
typed semantic edges, since graph compilation will likely become more expensive.

\subsection{4. Packing policy and schedule families}

The current row-compact and generalized 2-D packing framework is a reusable
``graph-to-kernel-family'' compiler. The main idea is already in place:
\begin{itemize}[leftmargin=1.5em]
    \item start from an irregular support graph,
    \item partition it into a small set of regular families,
    \item run specialized kernels on each family.
\end{itemize}

That is exactly the right systems pattern for richer semantic graphs too.

\section{What Does \emph{Not} Recycle Directly}

Not every part of the current HSA stack is directly reusable for a leap toward
semantic or causal graphs.

\subsection{1. Binary support is too weak for typed semantics}

If all edges are reduced to a single binary mask, then
\[
\texttt{Spain} \to \texttt{Europe}
\quad\text{and}\quad
\texttt{rain} \to \texttt{wet}
\]
look identical at the routing level. The graph topology survives, but relation
type does not.

\subsection{2. Standard softmax attention only models positive mixing weights}

Softmax produces
\[
\alpha_{ij} \ge 0, \qquad \sum_j \alpha_{ij} = 1.
\]
This is appropriate for compatibility-based readout, but it is not a natural
model of signed causal influence. The output can still become negative through
the value vectors, but the weighting mechanism itself is not signed.

\subsection{3. Current kernels assume binary support plus standard reductions}

The current sparse kernels assume:
\begin{itemize}[leftmargin=1.5em]
    \item allowed or forbidden edges,
    \item standard score computation,
    \item standard softmax normalization,
    \item standard $QK^\top$ and weighted $V$ reductions.
\end{itemize}

Typed, signed, or relation-conditioned message passing would require either:
\begin{itemize}[leftmargin=1.5em]
    \item multiple edge-family kernels,
    \item relation-specific score transforms,
    \item relation-specific value transforms,
    \item or an alternative aggregator that is no longer a pure masked softmax.
\end{itemize}

\section{From Hierarchy to Typed Directed Semantic Graphs}

\subsection{The right abstraction}

The natural extension is not ``arbitrary graph attention'' in the abstract. The
natural extension is:
\begin{quote}
decoder-causal sparse routing over tokens and semantic memory nodes, with typed
directed edges.
\end{quote}

This still respects autoregressive training, but allows the graph to encode
more than document boundaries.

\subsection{Semantic direction versus attention direction}

There are two distinct notions of direction:
\begin{enumerate}[leftmargin=1.5em]
    \item \textbf{semantic direction:} e.g.\ a knowledge or event relation,
    \item \textbf{attention-flow direction:} which previously visible node a
    token is allowed to read from.
\end{enumerate}

These are not the same thing. A semantic graph may contain a relation
\[
\texttt{rain} \to \texttt{wet},
\]
but in a decoder we still must causalize it into past-visible reads only. So a
semantic graph must be compiled into a causal routing graph before it becomes an
attention support.

\subsection{A minimal typed extension}

The smallest meaningful extension from current HSA is:
\begin{itemize}[leftmargin=1.5em]
    \item keep binary support families,
    \item add edge type $\tau_{ij}$,
    \item modify the score and/or transmitted message by type.
\end{itemize}

A minimal relation-aware score is:
\[
\mathrm{score}_{ij}
=
q_i^\top W_{\tau_{ij}} k_j + b_{\tau_{ij}}.
\]

A minimal relation-aware message is:
\[
m_{ij}
=
\alpha_{ij}\, T_{\tau_{ij}}(v_j),
\qquad
o_i = \sum_{j} m_{ij}.
\]

This already gives a typed directed semantic graph while remaining relatively
close to the current masked-attention implementation model.

\section{Toward Signed and Causal Message Passing}

\subsection{Why signed edges are interesting}

Some relations are naturally excitatory while others are inhibitory. For
example, a world-state graph might treat
\begin{itemize}[leftmargin=1.5em]
    \item \texttt{rain} $\to$ \texttt{wet-ground} as positive,
    \item \texttt{watering} $\to$ \texttt{wet-ground} as positive,
    \item \texttt{sun/heat} $\to$ \texttt{wet-ground} as negative or drying.
\end{itemize}

A pure binary mask cannot distinguish these. Even a typed score only changes who
is read from, not whether the message should reinforce or inhibit the target.

\subsection{What softmax cannot express directly}

Softmax gives positive simplex weights. So if the goal is to model signed or
inhibitory influence directly, one likely needs one of:
\begin{enumerate}[leftmargin=1.5em]
    \item relation-specific value transforms with sign,
    \item dual positive/negative channels,
    \item a non-simplex aggregator,
    \item latent semantic state nodes where inhibitory structure is carried in
    state rather than in the raw attention weights.
\end{enumerate}

\subsection{A realistic near-term precursor}

The most practical precursor to causal reasoning is not true interventional
causality. It is:
\begin{quote}
typed directed sparse message passing over tokens and semantic memory nodes,
where the graph encodes permitted information pathways and relation-specific
message transforms.
\end{quote}

That is already a meaningful leap beyond hierarchy, while remaining compatible
with autoregressive training.

\section{Why This Is Not Yet True Causal Reasoning}

A directed semantic graph inside a next-token predictor is still an
\emph{inductive bias}, not a causal model in Pearl's sense. True causal claims
require more than directional edges:
\begin{itemize}[leftmargin=1.5em]
    \item intervention semantics,
    \item counterfactual structure,
    \item assumptions about latent confounding,
    \item or explicit world-state modeling.
\end{itemize}

So the better description is:
\begin{quote}
the current HSA machinery can be extended into a typed semantic routing
framework that may support precursors to causal reasoning, but not causal
inference merely by virtue of using a directed graph.
\end{quote}

\section{A Concrete Leapfrog Plan}

\subsection{Stage 1: typed sparse graph}

Extend the current binary support into a small set of typed edge families:
\begin{itemize}[leftmargin=1.5em]
    \item local,
    \item summary/hierarchy,
    \item entity-link,
    \item discourse,
    \item causal-positive,
    \item causal-negative.
\end{itemize}

This step can reuse:
\begin{itemize}[leftmargin=1.5em]
    \item token-side metadata extraction,
    \item graph compilation,
    \item sidecars,
    \item packed family scheduling.
\end{itemize}

\subsection{Stage 2: relation-aware score and message transforms}

Keep the sparse routing framework, but make score and message depend on edge
type. This is still close enough to current HSA that the implementation can
evolve incrementally.

\subsection{Stage 3: semantic memory nodes}

Introduce explicit non-token nodes such as:
\begin{itemize}[leftmargin=1.5em]
    \item entity nodes,
    \item event nodes,
    \item state nodes,
    \item section or chapter summaries,
    \item retrieved knowledge nodes.
\end{itemize}

Then sparse attention becomes routing between tokens and memory/state nodes,
rather than only between raw tokens.

\subsection{Stage 4: signed or non-simplex aggregation}

Only after typed graph routing is stable should the scoring rule itself be
generalized beyond pure masked softmax, for example by allowing signed message
coefficients or separate positive/negative channels.

\section{Summary}

The current HSA implementation should be understood as \emph{sparse causal
routing}, not arbitrary graph attention. Its graph compiler, exact dense/sparse
references, sidecar system, and packing framework already provide reusable
infrastructure for a richer semantic-graph program.

The key engineering lessons so far are:
\begin{enumerate}[leftmargin=1.5em]
    \item correctness depends on semantic mask, sparse schedule, and kernel
    support all matching exactly;
    \item forward benefits strongly from aggressive packing because the problem
    is query-owned and easy to regularize;
    \item packed training still uses the exact sparse backward, but long-context
    runtime now depends critically on choosing the block-sparse traversal
    instead of the dense fallback;
    \item generalized 2-D packing is the right forward abstraction because it
    extracts dense submatrices from the support matrix itself;
    \item the natural next leap is typed directed semantic routing, not a jump
    straight to unrestricted signed graph inference.
\end{enumerate}

In short: the current method will not work with \emph{any} graph and token
scheme, but it is already close to a general graph compiler for \emph{causalized
sparse semantic routing}. That is a strong base for the next stage.

\section*{References}

\begin{thebibliography}{99}

\bibitem{vaswani2017attention}
Ashish Vaswani, Noam Shazeer, Niki Parmar, Jakob Uszkoreit, Llion Jones,
Aidan N. Gomez, Lukasz Kaiser, and Illia Polosukhin.
\newblock \emph{Attention Is All You Need}.
\newblock NeurIPS, 2017.
\newblock \url{https://arxiv.org/abs/1706.03762}

\bibitem{dao2022flashattention}
Tri Dao, Daniel Y. Fu, Stefano Ermon, Atri Rudra, and Christopher Re.
\newblock \emph{FlashAttention: Fast and Memory-Efficient Exact Attention with IO-Awareness}.
\newblock NeurIPS, 2022.
\newblock \url{https://arxiv.org/abs/2205.14135}

\bibitem{dao2023flashattention2}
Tri Dao.
\newblock \emph{FlashAttention-2: Faster Attention with Better Parallelism and Work Partitioning}.
\newblock 2023.
\newblock \url{https://arxiv.org/abs/2307.08691}

\bibitem{beltagy2020longformer}
Iz Beltagy, Matthew E. Peters, and Arman Cohan.
\newblock \emph{Longformer: The Long-Document Transformer}.
\newblock 2020.
\newblock \url{https://arxiv.org/abs/2004.05150}

\bibitem{zaheer2020bigbird}
Manzil Zaheer, Guru Guruganesh, Kumar Avinava Dubey, Joshua Ainslie,
Chris Alberti, Santiago Ontanon, Philip Pham, Anirudh Ravula,
Qifan Wang, Li Yang, and Amr Ahmed.
\newblock \emph{Big Bird: Transformers for Longer Sequences}.
\newblock NeurIPS, 2020.
\newblock \url{https://arxiv.org/abs/2007.14062}

\bibitem{velivckovic2018gat}
Petar Velickovic, Guillem Cucurull, Arantxa Casanova, Adriana Romero,
Pietro Lio, and Yoshua Bengio.
\newblock \emph{Graph Attention Networks}.
\newblock ICLR, 2018.
\newblock \url{https://arxiv.org/abs/1710.10903}

\bibitem{schlichtkrull2018rgcn}
Michael Schlichtkrull, Thomas N. Kipf, Peter Bloem, Rianne van den Berg,
Ivan Titov, and Max Welling.
\newblock \emph{Modeling Relational Data with Graph Convolutional Networks}.
\newblock ESWC, 2018.
\newblock \url{https://arxiv.org/abs/1703.06103}

\bibitem{battaglia2018relational}
Peter W. Battaglia, Jessica B. Hamrick, Victor Bapst, Alvaro Sanchez-Gonzalez,
Vinicius Zambaldi, Mateusz Malinowski, Andrea Tacchetti, David Raposo,
Adam Santoro, Ryan Faulkner, Caglar Gulcehre, H. Francis Song,
Andrew J. Ballard, Justin Gilmer, George E. Dahl, Ashish Vaswani,
Kelsey R. Allen, Charles Nash, Victoria Langston, Chris Dyer,
Nicolas Heess, Daan Wierstra, Pushmeet Kohli, Matthew Botvinick,
Oriol Vinyals, Yujia Li, and Razvan Pascanu.
\newblock \emph{Relational Inductive Biases, Deep Learning, and Graph Networks}.
\newblock 2018.
\newblock \url{https://arxiv.org/abs/1806.01261}

\bibitem{hu2020hgt}
Ziniu Hu, Yuxiao Dong, Kuansan Wang, and Yizhou Sun.
\newblock \emph{Heterogeneous Graph Transformer}.
\newblock WWW, 2020.
\newblock \url{https://arxiv.org/abs/2003.01332}

\bibitem{derr2018signedgcn}
Tyler Derr, Yao Ma, and Jiliang Tang.
\newblock \emph{Signed Graph Convolutional Network}.
\newblock ICDM, 2018.
\newblock \url{https://arxiv.org/abs/1808.06354}

\bibitem{pearl2009causality}
Judea Pearl.
\newblock \emph{Causality: Models, Reasoning, and Inference}.
\newblock Cambridge University Press, second edition, 2009.
\newblock \url{https://www.cambridge.org/core/books/causality/B0046844FAE10CBF274D4ACBDAEB5F5B}

\end{thebibliography}

\end{document}
